Se seleccionó el conjunto de datos CIC-IDS-2017 como base experimental del trabajo. Este conjunto de datos fue creado por el Canadian Institute for Cybersecurity (CIC) de la University of New Brunswick (UNB) y contiene un gran volumen de datos y variedad de ataques realistas y comunes. La infraestructura que utilizaron para la captura de datos incluía una gran variedad de dispositivos y sistemas operativos, con el objetivo de que fuera lo más realista posible. La captura se reparte en cinco días, de lunes a viernes, dividiendo los ataques en distintos días. Para generar tráfico humano de fondo, generaron tráfico benigno realista basado en la actividad real de 25 usuarios.

Una vez generado y capturado el tráfico de esos cinco días, utilizaron su herramienta CICFlowMeter para agrupar los paquetes en flujos bidireccionales. Cada flujo bidireccional representa una comunicación completa extremo-extremo, es decir, el tráfico intercambiado en dos extremos en ambos sentidos. Cada flujo está etiquetado como benigno o asociado a un ataque en específico y, además, tiene asociado un total de 83 características de tráfico que facilitarán el análisis y clasificación de estos. El material que proporcionan en su sitio web incluye ficheros csv en los que cada fila equivale a un flujo y cada columna a una característica y los PCAP originales capturados. Los flujos se agrupan en estos csv mediante la quíntupla IP origen-IP destino-Puerto origen-Puerto destino-Protocolo. Estas cinco son ejemplos de características que podemos encontrar, además de otras más complejas relacionadas con la duración del flujo, el número de paquetes, el volumen de bytes transferidos o los intervalos temporales entre paquetes.

En cuanto a los ataques que podremos clasificar en este conjunto de datos, se encuentran ataques de botnet, denegación de servicio distribuido (DDoS), distintas variantes de denegación de servicio (DoS), ataques de fuerza bruta sobre servicios FTP y SSH, explotación de la vulnerabilidad Heartbleed, actividades de infiltración, escaneo de puertos y ataques dirigidos a aplicaciones web, entre ellos inyección SQL (SQLi) y cross-site scripting (XSS).

En el material de partida se encuentran 5 de estos ficheros PCAP, uno por cada día. En cuanto a los ficheros CSV, hay algunos días para los que se tiene más de uno de estos. Esto es debido a que decidieron clasificar cada CSV con distintos tipos de ataques. Hay un fichero para el lunes, el cual contiene únicamente tráfico benigno; uno para el martes, que contiene ataques de fuerza bruta y tráfico benigno; uno para el miércoles, que contiene ataques de denegación de servicios; dos para el jueves, uno que contienen ataques web y otro que contiene ataques de infiltración; y tres para el viernes, uno con escaneos de puertos, otro con ataques de denegación de servicios y otro con ataques de botnet.